\section{Logarithmic--Stirling and supernecklace applications}
\label{sec:supernecklace}

The single-spike theorem above transfers an isolated local transposition to
the polylogarithmic and factorially weighted families in
\eqref{eq:transfer-families}.  The two scopes now diverge.  The strict Kaluza
argument in \Cref{thm:polylog-global} proves absolute indecomposability for
every pure polylogarithmic member and yields global symmetric families.  For
the factorially weighted $f_{\sigma,\tau}$ family the transfer remains local:
without a separate indecomposability theorem it does not by itself imply full
symmetric monodromy.  The logarithmic member $\operatorname{Li}_1$ has the
additional enumerative refinement developed below.

Put
\[
 C_m(X):=(m-1)!P_m^{(-\log(1-z))}(X)
 =\sum_{k=1}^m (k-1)!\stirling{m}{k}X^k,
 \tag{SN-1}\label{eq:SN-1}
\]
where \(\stirling{m}{k}\) is the unsigned Stirling number of the first kind.

\begin{proposition}[Type-III supernecklace enumerator; \UPASS]
The exponential generating function of the polynomials \(C_m\) is
\[
 \boxed{\displaystyle
 \sum_{m\ge1}C_m(X)\frac{z^m}{m!}
 =\log\!\left(\frac{1}{1-X\log\frac1{1-z}}\right).}
 \tag{SN-2}
\]
Consequently \(C_m\) is the cycle-number enumerator for a labelled cycle of
labelled cycles (the type-III supernecklace construction).
\end{proposition}

\begin{proof}
Expand the outer logarithm and use the classical Stirling exponential
generating function:
\[
 \log\!\left(\frac1{1-X L(z)}\right)
 =\sum_{k\ge1}\frac{X^k}{k}L(z)^k,
 \qquad
 L(z)^k=k!\sum_{m\ge k}\stirling{m}{k}\frac{z^m}{m!},
\]
where \(L(z)=\log(1/(1-z))\).  Coefficient comparison gives
\eqref{eq:SN-1}.  The labelled product/cycle interpretation is precisely the
standard exponential-formula reading of the same identity.
\end{proof}

\begin{theorem}[Symmetric supernecklace family; \UPASS]
For every prime \(p\ge5\), with \(m=p+2\),
\[
 \Gal(C_m(X)-T/\overline{\Q}(T))
 =\Gal(C_m(X)-T/\Q(T))=S_m.
 \tag{SN-3}\label{eq:SN-3}
\]
Moreover the off-diagonal collision polynomial
\[
 \boxed{\displaystyle
 \frac{C_m(X)-C_m(Y)}{X-Y}}
 \tag{SN-4}
\]
is absolutely irreducible over \(\Q\).
\end{theorem}

\begin{proof}
Scaling a polynomial by the nonzero constant \((m-1)!\) does not change its
geometric or arithmetic monodromy.  Equation~\eqref{eq:stirling-mon} therefore
gives \eqref{eq:SN-3}.  For a polynomial cover, the irreducible factors of the
off-diagonal fiber product correspond to the orbits of a point stabilizer on
the remaining sheets.  The natural action of \(S_m\) is two-transitive, so
that stabilizer has one such orbit.  This proves absolute irreducibility.
\end{proof}

\begin{remark}[Boundary of the application]
Qualitative Hilbert irreducibility supplies many specializations with group
\(S_m\), but no quantitative exceptional-set exponent is asserted here.
Nor is a genus formula claimed for the collision curve: symmetric
monodromy gives absolute irreducibility, while smoothness and the precise
ramification divisor require a separate Morse analysis.
\end{remark}
